\section{Dataset and Preprocessing}
\label{sec:dataset}

MVTec AD~\cite{bergmann2019mvtec,bergmann2021mvtec} contains 15 categories.
Each training set contains only defect-free images, and each test set mixes good images with defects annotated by pixel-precise ground-truth (GT) masks.
We use three objects (bottle, capsule, screw) and three textures (carpet, tile, leather); Table~\ref{tab:dataset} summarises them.
Using a fixed seed, we split each \texttt{train/good} set into 90\% for fitting and 10\% (21--32 images) for a \emph{normal-only validation set}, which is never used for fitting.

\begin{table}[t]
  \centering
  \caption{Selected categories. Fit/val: our split of train/good. Area: median GT defect area (\% of image). Pix.\ std: mean per-pixel standard deviation of normal images. Align: mean correlation of a normal image with the mean normal image.}
  \label{tab:dataset}
  \footnotesize
  \setlength{\tabcolsep}{2.5pt}
  \begin{tabular}{@{}llrrrrrr@{}}
\toprule
Category & Kind & Fit/val & \shortstack{Test\\G/D} & Types & \shortstack{Area\\\%} & \shortstack{Pix.\\std} & Align \\
\midrule
bottle & obj. & 188/21 & 20/63 & 3 & 5.72 & 0.022 & 0.996 \\
capsule & obj. & 197/22 & 23/109 & 5 & 0.43 & 0.023 & 0.998 \\
screw & obj. & 288/32 & 41/119 & 5 & 0.29 & 0.094 & 0.579 \\
carpet & tex. & 252/28 & 28/89 & 5 & 1.48 & 0.058 & 0.200 \\
tile & tex. & 207/23 & 33/84 & 5 & 7.93 & 0.106 & 0.155 \\
leather & tex. & 220/25 & 32/92 & 5 & 0.46 & 0.022 & 0.880 \\
\bottomrule
\end{tabular}

\end{table}

\begin{figure}[t]
  \centering
  \includegraphics[width=\columnwidth]{fig_defect_area.pdf}
  \caption{Size of every GT defect region inside the 224-pixel evaluation crop, in stride-8 feature cells (log scale; one cell $=8\times8$\,px $=0.128$\% of the crop). Each point is one region (jittered vertically); boxes show the quartiles and whiskers extend to 1.5 times the interquartile range, both computed on the log scale; $n$ is the number of regions. The dashed line marks one feature cell, and the faint lines mark the size bins of Fig.~\ref{fig:regionsize}. Many screw, capsule and leather defects cover only a few cells.}
  \label{fig:defectarea}
\end{figure}

Four dataset characteristics shape the difficulty of each category.
\textbf{Defect size} spans about three orders of magnitude, from 0.037\% of the image (the smallest capsule crack) to 35.49\% (a tile glue strip). Fig.~\ref{fig:defectarea} shows the size of every GT region in the units the models see.
Screw, capsule and leather have median defect areas below 0.5\%.
At a 224-pixel input, one stride-8 WRN50 feature cell covers about 0.13\% of the image, so many of these defects span only a few cells.
\textbf{Alignment} is near-perfect for bottle and capsule (align $\approx$\,0.996--0.998), while screws lie at arbitrary rotation and position (0.579).
This is the property that breaks position-specific models.
\textbf{Texture versus object}: textures are stationary, so any patch can be compared with any normal patch and alignment does not matter (for textures, \emph{align} reflects only illumination).
For objects, the decision rests on a small foreground region against a fixed background, where specks and illumination gradients can trigger false alarms.
\textbf{Normal variability} is highest for tile and screw (pix\_std 0.106 and 0.094).
High variability broadens the learned normal distribution and makes threshold calibration harder.
Finally, up to 10\% of anomalous images contain several separate defect regions, and AU-PRO weights each region equally.

\textbf{Preprocessing.}
Images are resized to 256 pixels, centre-cropped to 224 and normalised with ImageNet mean and standard deviation.
Masks receive the same geometric transform, and all metrics are computed at 224$\times$224.
Matching the backbone's pre-training statistics is essential for frozen features.
The centre crop follows the reference implementations but discards a 12.5\% border, which removes part of a few defects (Section~\ref{sec:failures}).
